\section{Auditoría de fuentes y mapa del curso: la misma attention, dos ejes de cómputo}

Esta clase se compone de dos partes que encajan entre sí. Primero, Aidan Gomez dedica unos quince minutos a repasar el origen del Transformer y la intuición detrás de su entrenamiento; después, Jannik Kossen y Neil Band presentan el Non-Parametric Transformer (NPT). Si solo se resumiera en el orden de la exposición, ambas partes parecerían «un repaso de fundamentos más un artículo»; sin embargo, el hilo didáctico real es más unitario: el Transformer coloca primero la attention \textbf{entre los tokens de una misma secuencia}, y NPT eleva después el eje sobre el que actúa la attention al plano \textbf{entre los datapoints de un mismo conjunto de datos}.

Esta nueva redacción toma como referencia temporal la publicación oficial de Stanford Online `zejXBg-2Vpk`. El video ofrece 1,476 subtítulos oficiales en inglés, hechos a mano y analizables, pero no hay un PDF de slides publicado por separado; por ello, se recuperaron de la grabación en 1080p y se seleccionaron manualmente 28 páginas didácticas. De las animaciones progresivas solo se conserva el estado final completo o los estados intermedios que tienen una función didáctica propia; se omiten las páginas anteriores a las que se regresa varias veces durante la sesión de Q\&A, las transiciones en blanco y el cierre del video, mientras que las explicaciones orales se incorporan al teacher-voice ledger y al cuerpo del texto.

\lecturefigure{slide-01-transformers-title.jpg}{Título de la primera parte: breve historia e intuición del Transformer.}{Video oficial de Stanford Online 00:00:46--00:01:07.}

Esta página de título recuerda que la tarea de Gomez no es volver a impartir un curso completo sobre el Transformer, sino extraer los tres mecanismos que determinan el resto de la historia: self-attention, multi-head attention y el autoregressive decoder que puede entrenarse en paralelo. Una vez comprendidos estos tres mecanismos, resulta más fácil ver que la innovación central de NPT no es «inventar otro tipo de attention», sino redefinir sobre qué objetos opera la attention.

\lecturefigure{slide-02-transformer-overview.jpg}{Panorama del Transformer: multi-head attention, self-attention y fast autoregressive decoding.}{Video oficial de Stanford Online 00:01:08--00:01:41.}

\begin{importantbox}{Una pregunta que condensa toda la clase}
Cuando el modelo predice un ejemplo, ¿solo puede apoyarse en «las características de ese ejemplo + la experiencia de entrenamiento ya comprimida en los parámetros», o puede leer directamente otros ejemplos de entrenamiento durante la inferencia? Una supervised neural network estándar suele adoptar la primera opción; NPT intenta aprender la segunda, y deja que la attention decida de extremo a extremo «qué ejemplos leer y cómo combinarlos».
\end{importantbox}

\teachervoice{Gomez resume la contribución del Transformer como tres componentes combinados, no como una fórmula aislada. Esta perspectiva es importante: LayerNorm, el learning-rate warmup, la inicialización y los detalles de las conexiones residual, que después se volvieron configuración estándar, eran en ese momento variables experimentales que podían determinar de forma notable el éxito o el fracaso.}

\subsection{Ruta de lectura}

A continuación se establece primero el marco matemático mínimo de la attention a nivel de token y del entrenamiento autoregresivo, y luego se pasa al input contract a nivel de conjunto de datos de NPT. La parte experimental no se limitará a repetir que «obtiene una muy buena posición en la clasificación», sino que responderá en orden a cuatro preguntas más exigentes: si NPT es suficientemente competitivo, si puede representar algoritmos que operan entre ejemplos, si de verdad utiliza otros ejemplos y qué estructura tiene esa dependencia.

\begin{knowledgebox}{No confundir los dos ejes}
\begin{tabularx}{\textwidth}{L{0.23\textwidth}>{\raggedright\arraybackslash}X >{\raggedright\arraybackslash}X}
\toprule
Mecanismo & Objetos que compara la attention & Pregunta principal \\
\midrule
Self-attention sobre secuencias & Los tokens / attributes dentro de un mismo ejemplo & ¿Qué elementos de la oración o del ejemplo están relacionados? \\
ABD de NPT & Los datapoints / rows de un mismo lote & ¿De qué otros ejemplos debe leer información el ejemplo actual? \\
ABA de NPT & Los attributes dentro de un mismo datapoint & ¿Cómo deben transformarse e interactuar los campos dentro de una fila? \\
\bottomrule
\end{tabularx}
\end{knowledgebox}

\subsection{Resumen de la sección}

Esta clase no es la unión de dos temas sin relación, sino la ampliación progresiva del ámbito sobre el que actúa la attention. La primera parte explica el entrenamiento en paralelo y las relaciones dentro de la secuencia; la segunda convierte el propio conjunto de datos en entrada del modelo, de modo que la predicción pueda depender explícitamente de otros datapoints.
